\Needspace{22\baselineskip}\SourceLeaf{31}
\[
\begin{array}{r|rrr}
n & a & b & c \\
\hline
2 & -1 & 1 & 0 \\
3 & 2 & 0 & -1 \\
4 & 1 & 1 & -1 \\
5 & 1 & -4 & 2 \\
6 & 1 & -6 & 3 \\
7 & 1 & -1 & 0 \\
8 & -2 & 0 & 1 \\
9 & -2 & 0 & 1 \\
10 & 1 & 6 & -3 \\
11 & 1 & 4 & -2 \\
12 & -107 & 33 & 6 \\
13 & 4 & 3 & -2 \\
14 & 1 & 2 & -1 \\
15 & 1 & -30 & 12 \\
16 & 1 & 50 & -20 \\
17 & 18 & -7 & 0 \\
19 & -8 & 3 & 0 \\
37 & 10 & -3 & 0 \\
\end{array}
\]

\Needspace{12\baselineskip}
\begingroup\small

\begin{center}\textbf{BEMERKUNGEN ÜBER DIE FRAGMENTE ZUR THEORIE DER AUS EINER CUBIKWURZEL
ZU BILDENDEN GANZEN ALGEBRAISCHEN ZAHLEN.}\end{center}\par\noindent\hspace*{\parindent}Die hier reproducirten Entwicklungen stammen ungefähr aus derselben Zeit wie
die voraufgehenden »Zur Theorie der cubischen Reste«. Nach einer
Tagebuchnotiz hat \textsc{Gauss} im speciellen die Untersuchung über Auflösung der
Gleichung $x^3+ny^3+nnz^3-3nxyz=1$ am 23. Dec. 1808 begonnen. Es liegt damit
ein neues Zeugniss für das so frühzeitige Auftreten ganzer algebraischer
Zahlen vor.

Im Sinne der neueren Zahlentheorie handelt es sich hier um ganze Zahlen
desjenigen reellen cubischen Körpers, der durch die reine Gleichung $x^3=n$
definirt ist. Es bedeutet $t$ eine ganze Zahl dieses Körpers, $\varphi t$ die
Norm von $t$ und $Ft$ das gleichfalls im Körper enthaltene Product der mit
$t$ conjugirten Zahlen, welche etwa $t_1$, $t_2$ heissen mögen:

\[
\varphi t = t t_1 t_2, \qquad Ft = t_1 t_2.
\]
